%================================================================
\section{What the procedure is not, and its limits}\label{sec:limits}

The procedure is easy to over-read, so we list what it does not claim.

\paragraph{It is not a posterior or a confidence set} NROY rules points out; it
does not weight the points that remain, and $\Theta_{\hm}$ is a range, not a
distribution of effects. The threshold $\tau$ is the $\Confidence$ quantile of a
$\chi^2_{\NMoments}$ distribution, but $V$ is diagonal while the two level moments
share a state stream, and the discrepancy $\delta$ rather than sampling noise
fills most of $V$. So $\tau$ is a tolerance, not a test level, and the image
carries no coverage guarantee. Approaches that do weight the parameter space, such
as Bayesian calibration with an explicit discrepancy term \citep{kennedy2001},
answer a different question.

\paragraph{It is not a coverage claim for the one-panel intervals} The two flags
describe one draw. The only coverage evidence is the Foil~2 calibration check
(Appendix~\ref{app:foil2}) and the repeated panels of Section~\ref{sec:a2results},
where the A1 and A2 rates reflect conservative intervals rather than calibration.

\paragraph{It does not make A1 or A2 correct} A1 and A2 are biased downward at
every point; transport shrinks the bias and does not remove it. Neither identifies
$\theta_2$ from the pilot, and their intervals target $\theta_2$ under assumptions
that the kernel does not satisfy exactly. They do not estimate the same functional
as A0, which equals $\theta_2$ only under no interference.

\paragraph{It is not a criticism of standard practice} A regression with robust
standard errors quantifies uncertainty within one panel under the analyst's stated
assumptions, and it does that well. An exposure-mapping estimator does the same
under its own assumptions \citep{aronow2017,leung2020}, and recent work studies
what such estimators recover when the map is wrong \citep{savje2021} and how to
test hypotheses about interference directly \citep{athey2018}. The overlay asks a
different question: whether one parameter point, or one world, is enough for a
policy report. The failures we document come from applying a one-world tool to a
question about the set of worlds the targets cannot rule out.

\paragraph{The analysts' panel overlaps the targets} The analysts' one panel is
stream offset $0$, which is also one of the $\RepsZ$ replications whose mean
defines $z$. The analysts therefore see data that helped set the targets, with a
weight of $1/\RepsZ$. We did not rerun the pipeline without it, and we have not
listed the rejected wave points just above $\tau$, which a small shift in $z$
could move into NROY.

\paragraph{It is one planted village} The project has one graph, one planted point
and one set of targets, and the targets are the simulator's own output. It makes
no claim about any real campaign and claims no external validity. With external
targets, $\delta$ becomes the allowance for a model that is not the world and must
be defended; the targets would also need a stated population and a model of how
they transport to the village of interest.

\paragraph{Who might use it} The composition is useful to anyone who reports a
policy number from a calibrated agent-based, structural or microsimulation model.
The one-point Monte Carlo error bar is a statement about one world. The image is a
statement about the worlds the targets cannot rule out, and its width is measured
rather than assumed. An empty NROY is also a usable result: it rejects the pair of
kernel and targets for the declared design. Where interference matters and the
data come from a randomized pilot, the ladder of Section~\ref{sec:analysts} makes
explicit what each step from the pilot to the policy assumes, and the planted
setting makes its cost visible. A design-based answer that relies on fewer of
these assumptions needs randomization at the level of the policy itself, for
instance across villages \citep{hudgens2008,baird2018}.
